\documentclass{article}
\usepackage[UTF8]{ctex}
\usepackage{amsmath,amssymb,amsfonts}
\usepackage{enumitem}
\usepackage{geometry}
\geometry{a4paper, margin=1in}
% === 错题本排版优化（不参与题目标识） ===
\setlength{\parindent}{0pt}
\setlength{\parskip}{0.35em plus 0.10em minus 0.08em}
\setlength{\emergencystretch}{3em}
\setlist{itemsep=0.25em, topsep=0.40em, parsep=0pt}
\allowdisplaybreaks[2]
% === 错题本排版优化结束 ===
\begin{document}
\section*{湖南大学2025年高等代数真题}
\section*{133．湖南大学2025年高等代数试卷（回忆版）}

\subsection{题目1}
1．设非零多项式 \(f(x), g(x) \in K[x], k\) 为正整数，满足

\[
(k+1) \operatorname{deg}(g(x))>\operatorname{deg}(f(x)) \geq k \operatorname{deg}(g(x)) .
\]

证明：存在 \(p_{0}(x), p_{1}(x), \cdots, p_{k}(x) \in K[x]\) ，满足 \(\operatorname{deg}\left(p_{i}(x)\right)<\operatorname{deg}(g(x))\) 或 \(p_{i}(x)=0\) ，使得

\[
f(x)=p_{0}(x)+p_{1}(x) g(x)+\cdots+p_{k}(x) g^{k}(x) .
\]

\subsection{题目2}
2．求解线性方程组

\[
\left\{\begin{array}{c}
x_{1}+2 x_{2}+\cdots+n x_{n}=n+1 ; \\
x_{1}+2^{2} x_{2}+\cdots+n^{2} x_{n}=(n+1)^{2} \\
\cdots \cdots \\
x_{1}+2^{n} x_{2}+\cdots+n^{n} x_{n}=(n+1)^{n}
\end{array}\right.
\]

\subsection{题目3}
3．给定向量组

\[
\alpha_{1}=(*, *, *, *, *), \alpha_{2}=(*, *, *, *, *), \alpha_{3}=(*, *, *, *, *), \beta_{1}=(*, *, *, *, *), \beta_{2}=(*, *, *, *, *) .
\]

记 \(V_{1}=L\left(\alpha_{1}, \alpha_{2}, \alpha_{3}\right), V_{2}=L\left(\beta_{1}, \beta_{2}\right)\) ，求 \(V_{1}+V_{2}\) 与 \(V_{1} \cap V_{2}\) 的一组基和维数。

\subsection{题目4}
4．对于 \(n\) 阶实对称矩阵 \(A\) ，存在 \(Y \in \mathbb{R}^{n}\) 使得 \(Y^{\mathrm{T}} A Y>0\) 。证明：对任意的 \(X \in \mathbb{R}^{n}\) ，有

\[
\left(X^{\mathrm{T}} A Y\right)^{2} \geq\left(X^{\mathrm{T}} A X\right)\left(Y^{\mathrm{T}} A Y\right)
\]

的充要条件是 \(X^{\mathrm{T}} A X\) 的正惯性指数等于 1 。

\subsection{题目5}
5．设 \(V=\mathbb{R}^{2 \times 2}, V\) 上的线性变换 \(\varphi\) 满足 \(\varphi(X)=A^{\mathrm{T}} X A, X \in V\) ，其中

\[
A=\left(\begin{array}{ll}
1 & 1 \\
0 & 1
\end{array}\right)
\]

求 \(V\) 中的一组基，使得 \(\varphi\) 在这组基下的矩阵为若尔当形矩阵。

\subsection{题目6}
6．设 \(A\) 是 \(n\) 阶实矩阵，且 \(\operatorname{det}(A)=0\) ，证明：\(r(A)=n-1\) 当且仅当存在列向量 \(X, Y \in \mathbb{R}^{n}\) ，使得

\[
\operatorname{det}\left(A+X Y^{\mathrm{T}}\right) \neq 0 .
\]

\subsection{题目7}
7．设 \(A, B\) 均为数域 \(K\) 上的 \(n\) 阶方阵。若存在数域 \(K\) 上的 \(n\) 阶可逆矩阵 \(P\) ，使得 \(P^{-1} A P\) 为上三角矩阵，则称 \(A\) 在数域 \(K\) 上可相似上三角化。证明：

\begin{itemize}
  \item[（1）] 数域 \(K\) 上的矩阵 \(A\) 可相似上三角化的充要条件是 \(A\) 的所有复特征值都在数域 \(K\) 中。
  \item[（2）] 若数域 \(K\) 上的矩阵 \(A, B\) 均可相似上三角化，且 \(A B=B A\) ，则 \(A, B\) 可同时相似上三角化，即存在可逆矩阵 \(Q\) ，使得 \(Q^{-1} A Q, Q^{-1} B Q\) 同时为上三角矩阵。
  \item[（3）] 若 \(A B=B A\) ，其中 \(B\) 还是幂零矩阵，证明： \(\operatorname{det}(A)=\operatorname{det}(A+B)\) ．
\end{itemize}

\subsection{题目8}
8．实矩阵 \(A=\left(a_{i j}\right)_{n \times n}\) 满足 \(a_{i i}>0, a_{i j}=-a_{j i}(i \neq j)\) ，证明： \(\operatorname{det}(A)>0\) ．

\subsection{题目9}
9．设 \(\mathcal{A}\) 为数域 \(F\) 上 \(n\) 维线性空间 \(V\) 上的线性变换，令

\[
V_{1}=\left\{\alpha \in V \mid \text { 存在正整数 } r \text {, 使得 } \mathcal{A}^{r}(\alpha)=0\right\}, V_{2}=\bigcap_{i=1}^{\infty} \mathcal{A}^{i}(V) \text {. }
\]

证明：

\begin{itemize}
  \item[（1）] \(V_{1}\) 和 \(V_{2}\) 都是 \(\mathcal{A}\) 的不变子空间。
  \item[（2）] \(\mathcal{A}\) 限制在 \(V_{1}\) 上是幂零变换．
  \item[（3）] \(\mathcal{A}\) 限制在 \(V_{2}\) 上是可逆变换．
  \item[（4）] \(V=V_{1} \oplus V_{2}\) ．
\end{itemize}
\end{document}
